\section{Graphs of geometries and the boundary of iteration}\label{sec:app-graphs}
A particular fundamental-group presentation can be compared by explicit changes of basis. This does not prove that arbitrary geometric HNN steps remain admissible in succession. The source's three-tree example has the exact positive comparison below; the subsequent two-loop example isolates the obstruction to an unrestricted iterative recipe.

\subsection{The three spanning trees in Example 5.18}
Use the six source generators and stable letters $t_G,t_H,t_K$.
For the first tree the relations are
\[
 a_1=b_4,\quad b_3=c_6,\quad t_G^{-1}c_5t_G=a_2,
 \quad t_G^{-1}c_6t_G=a_1.
\]
Eliminate four names, obtaining the free basis
$x=b_3=c_6$, $y=c_5$, $t=t_G$, with
$a_1=b_4=t^{-1}xt$ and $a_2=t^{-1}yt$.
For the second tree use $x=a_1=b_4=c_6$, $y=a_2=c_5$, $t=t_H$,
and eliminate $b_3=t^{-1}xt$. For the third use
$x=a_1=c_6=b_3$, $y=a_2=c_5$, $t=t_K$,
and eliminate $b_4=t^{-1}xt$.
Every substitution is reversible and leaves no relator on $x,y,t$;
these are explicit isomorphisms onto $F(x,y,t)$.

For each tree the final types are $X=\{1,3,4,6\}$,
$Y=\{2,5\}$ and the stable type. Their maximal parabolics become
\[
 \langle y,t\rangle,\qquad\langle x,t\rangle,\qquad\langle x,y\rangle.
\]
For instance the first tree's $X$ subgroup is
$\langle a_2,c_5,t_G\rangle=\langle t^{-1}yt,y,t\rangle
=\langle y,t\rangle$, and its stable subgroup is
$\langle b_3,c_5,c_6\rangle=\langle x,y\rangle$.
All others follow from the displayed substitutions. Thus the same
group isomorphisms identify the actual parabolics and induce coset
incidence isomorphisms, establishing the full example. Each is the
standard rank-three free-group geometry, with infinite panels.
On source page 56, $a_2=c_5$ is incorrectly claimed also for the
first tree; there they are conjugate. Merely finding separately
isomorphic parabolic groups would not suffice for an incidence
isomorphism, whereas the explicit marked maps do.

This example does not settle the general fundamental-geometry or
universal-property questions in Sections 5.4--6. The obstruction below concerns only the naive sequential algorithm. All source
application attributions remain with Piedade and Tranchida; no
historical-priority or new-question claim is made here.



\subsection{Failure of sequential admissibility}
Piedade--Tranchida, arXiv:2505.24662v1, Section 5.4, proposes graphs of
coset incidence systems and asks whether the amalgamation and HNN
operations give a fundamental geometry independent of the spanning tree.
The following example isolates an earlier obstruction: locally admissible
edge maps need not remain admissible after the first HNN operation.
It concerns the stated iterative procedure, not the impossibility of
defining a different fundamental geometry.

\begin{proposition}
There is a graph with one vertex and two loops, all of whose vertex and
edge geometries are the same finite regular hypertope and all of whose
edge identifications are admissible, for which neither ordering of the
two prescribed geometric HNN operations is admissible throughout.
\end{proposition}
\begin{proof}
Let $A=C_2^3$ with basis $e_1,e_2,e_3$, and give it the standard coset
geometry with maximal parabolics
\[
 A_i=\langle e_j:j\ne i\rangle\quad(i=1,2,3).
\]
This is the join of three two-point geometries, hence is thin,
residually connected and flag-transitive, with trivial Borel.
The graph has this geometry at its sole vertex and at both loops.
On each loop take one boundary map to be the identity. Take the other
to be respectively the coordinate transpositions
$\sigma=(12)$ and $\tau=(23)$. These are actual automorphisms
permuting every indexed parabolic, so each edge map is admissible.

Perform the $\sigma$ operation first and call its stable letter $t$.
The resulting group is $G=A\rtimes_\sigma\langle t\rangle$, with
$t$ of infinite order. The type classes and maximal parabolics prescribed
by the source construction are
\[
 \{1,2\},\quad\{3\},\quad *,\qquad
 G_{\{1,2\}}=\langle e_3,t\rangle,
 \quad G_{\{3\}}=\langle e_1,e_2,t\rangle,
 \quad G_*=A.
\]
The standard parabolics contained in $A$ are exactly
\[
 A,\qquad \langle e_3\rangle,\qquad
 \langle e_1,e_2\rangle,\qquad 1.
\]
Indeed every intersection not involving $G_*$ contains the infinite-order
element $t$, and the intersections involving it are precisely the four
groups displayed. The second edge still identifies the parabolic $A$
with itself by $\tau$, but
\[
 \tau(\langle e_3\rangle)=\langle e_2\rangle
\]
is not on this list. Thus $\tau$ is not admissible for the new geometry.

If the $\tau$ operation is performed first, its two old type classes
are $\{1\}$ and $\{2,3\}$. The proper nontrivial parabolics in its
$A$ residue are $\langle e_1\rangle$ and $\langle e_2,e_3\rangle$.
The remaining map $\sigma$ sends $\langle e_1\rangle$ to
$\langle e_2\rangle$, again absent. Hence the reverse ordering also
fails. Both transpositions are their own inverses, so reversing loop
orientation does not repair either order.
\end{proof}

The abstract fundamental group still exists: it has presentation
\[
 \langle A,t,u\mid t^{-1}at=\sigma(a),\ u^{-1}au=\tau(a)
 \ (a\in A)\rangle
 \cong A\rtimes F(t,u).
\]
The obstruction is in the parabolic geometry retained after one operation.
A simultaneous construction, a common refinement of the retained incidence
data, or a global choice of type classes might avoid it; none is supplied
by this example. The source question must therefore distinguish an
unrestricted iterative recipe from a construction with additional coherent
type data. This limits the proposed iterative recipe; it is not an answer to the broader existence question.
